\documentclass{article}

\usepackage[final]{project_nips_style}

% =========================
% Packages
% =========================
\usepackage[utf8]{inputenc}
\usepackage[T1]{fontenc}
\usepackage{hyperref}
\usepackage{url}
\usepackage{booktabs}
\usepackage{amsmath, amsfonts}
\usepackage{graphicx}
\usepackage{microtype}
\usepackage{xcolor}


% =========================
% Title Information
% =========================
\title{Multimodal Skin Lesion Classification: Investigating the Influence of Metadata on Model Decision-Making}

\author{
  Eva Bauer \\
  University of Bern \\
  \texttt{eva.bauer@students.unibe.ch}
  \And
   Alexandra Benova \\
  University of Bern \\
  \texttt{alexandra.benova@students.unibe.ch}
   \And
   Nina Fröhlich \\
  University of Bern \\
  \texttt{nina.froehlich@students.unibe.ch}
}

\begin{document}

\maketitle


\vspace{0.5em}


% ============================================================
\section{Introduction and Background}
% ============================================================

Skin cancer is one of the most common cancer types around the world, which makes exploring it so crucial,
as also early detection makes a big difference in treatment chances \cite{watson2026}.
A more recent approach to image-based prediction models is the use of multimodal data, which provides a more detailed
view on the samples and improves diagnosis and treatment \cite{hao2025}. In oncology, multimodal approaches enhance tumor
characterization and enable more personalized treatment planning \cite{hao2025}. 
Here we investigated how patient metadata influences the model decision-making specifically in skin lesion
classification. We compared image-only, metadata-only and multimodal models and compared their performance.
We analyzed which metadata features contribute most to predictions and how they alter visual attention.
To do so, we used advanced explainability techniques (Grad-CAM and CRP) to visualize the model attention and
feature relevance.


% ============================================================
\section{Method}
% ============================================================
\subsection{Model Architecture}

Our model combines dermoscopic images and patient metadata using a multimodal architecture consisting of three main components: an image encoder, a metadata encoder, and a fusion module.

For the image encoder, we used a convolutional neural network (CNN) (ResNet \cite{he2016resnet}), since CNNs perform well on smaller datasets and are effective at capturing local texture patterns in skin lesions \cite{esteva2017dermatologist}. The input image is mapped to a latent feature vector representing relevant visual characteristics such as asymmetry, border irregularity, and color variation.

The metadata (e.g. age, sex, anatomical site, skin type) is processed using a TabTransformer \cite{huang2020tabtransformer}. Each feature is treated as an individual token and embedded into a vector space. Self-attention layers are then applied to model relationships between features, allowing the model to capture dependencies such as how age interacts with lesion location.

To combine both modalities, we use a cross-attention mechanism instead of simple concatenation. This allows the image features to dynamically attend to relevant metadata. Following the standard attention formulation \cite{vaswani2017attention}, the image embedding is used as the query (Q), while metadata embeddings are used as keys (K) and values (V). This results in a weighted combination of metadata features conditioned on the image. The fused representation is then passed to a small MLP for binary classification (malignant vs. benign).

An advantage of this approach is its inherent explainability. The attention weights can be extracted to quantify how much each metadata feature contributes to a specific prediction. In addition, we apply Grad-CAM \cite{selvaraju2017gradcam} to visualize the image regions that are most relevant for the model decision.

The model is trained using binary cross-entropy loss and the Adam optimizer \cite{kingma2014adam}. To improve generalization, we apply standard data augmentation and dropout.


\subsection{Training Setup}

\textbf{Dataset} \
For our analysis we used the Melanoma Research Alliance Multimodal Image Dataset (MRA-MIDAS) \cite{midas2023}. It represents the first publicly available dataset with prospectively collected and systematically paired dermoscopic and clinical images covering multiple skin lesion types \cite{midas2023}. The dataset includes 3216 images and metadata from 733 patients and is relatively balanced between benign (1500) and malignant (1419) cases.

The metadata contains patient-level information such as age, ethnicity, and sex, as well as lesion-specific attributes including size, anatomical location, and textual descriptions. Ground-truth labels are provided via “midas\_path”, and up to three additional clinical impressions are available per sample.

We performed binary classification into malignant (=1) and benign (=0) based on the “midas\_path” labels. The labeled data was split into training, validation, and test sets (80\%-10\%-10\%). Additionally, 497 samples without ground-truth labels were used as an external test set, where ‘clinical\_impression\_1’ served as a proxy label.


\subsection{Implementation Details}

All models were initialized with ImageNet-pretrained weights for the image encoder. Images were resized and augmented using standard techniques such as random flipping, rotation, and color jitter to improve generalization.

For the metadata, categorical features were embedded and numerical features were normalized before being passed to the TabTransformer. During training, dropout was applied to reduce overfitting.

We used the Adam optimizer with a fixed learning rate and trained the models for a predefined number of epochs, selecting the best model based on validation performance.


% ============================================================
\section{Results}
% ============================================================




% ============================================================
\section{Conclusion}
% ============================================================

\newpage
{\small
\bibliographystyle{plainnat}
\bibliography{references}
}

% ============================================================
\appendix
% ============================================================

\newpage
\section{Appendix (Optional)}



\end{document}